\PassOptionsToPackage{dvipsnames,svgnames,table}{xcolor}
\documentclass[11pt]{article}
\usepackage[a4paper,margin=22mm]{geometry}
\usepackage[no-math]{fontspec}
\setmainfont{DejaVu Serif}
\usepackage{amsmath,amssymb,amsthm,mathtools,graphicx,xcolor,hyperref,xurl,xspace,ifthen,enumerate}
\usepackage[shortlabels]{enumitem}
\setlength{\emergencystretch}{4em}
\AtBeginDocument{\special{pdf:mapline pzdr Dingbats <uzdr.pfb}}
\SetMathAlphabet{\mathbf}{normal}{OT1}{cmr}{bx}{n}
\SetMathAlphabet{\mathbf}{bold}{OT1}{cmr}{bx}{n}
\newtheorem{theo}{Theorem}[section]
\newtheorem{lemm}[theo]{Lemma}
\newtheorem{thm}[theo]{Theorem}
\newtheorem{proposition}[theo]{Proposition}
\newtheorem{corollary}[theo]{Corollary}
\newtheorem{definition}[theo]{Definition}
\newtheorem{theorem}[theo]{Theorem}
\newtheorem*{remas*}{Remarks}
\newtheorem{remas}[theo]{Remarks}
\newtheorem*{exam*}{Example}
\newtheorem{ques}[theo]{Question}
\providecommand{\og}{«\nobreakspace}
\providecommand{\fg}{\nobreakspace»}
\newtheorem{lem}[theo]{Lemma}
\newtheorem{lemma}[theo]{Lemma}
\newtheorem{prop}[theo]{Proposition}
\newtheorem{coro}[theo]{Corollary}
\newtheorem{corol}[theo]{Corollary}
\newtheorem*{conj*}{Conjecture}
\newtheorem{conj}[theo]{Conjecture}
\newtheorem{conjecture}[theo]{Conjecture}
\newtheorem{Proposition}[theo]{Proposition}
\newtheorem{theoreme}[theo]{Theorem}
\theoremstyle{definition}
\newtheorem{defi}[theo]{Definition}
\newtheorem{exam}[theo]{Example}
\newtheorem{exem}[theo]{Example}
\newtheorem{remark}[theo]{Remark}
\newtheorem{example}[theo]{Example}
\newtheorem{exams}[theo]{Examples}
\newtheorem{rema}[theo]{Remark}
\newtheorem{nota}[theo]{Notation}
\newtheorem{notation}[theo]{Notation}
\newtheorem*{notas}{Notations}
\newtheorem{result}[theo]{Result}
\newtheorem{question}[theo]{Question}
\newtheorem{prob}[theo]{Problem}
\newtheorem*{rema*}{Remark}
\newtheorem*{defi*}{Definition}
\newtheorem*{theo*}{Theorem}
\newtheorem*{lemm*}{Lemma}
\newtheorem*{prop*}{Proposition}
\newcommand{\enoncename}{}
\newtheorem{corpusenonce}[theo]{\enoncename}
\newenvironment{enonce}[2][]{\renewcommand{\enoncename}{#2}\begin{corpusenonce}}{\end{corpusenonce}}
\newenvironment{enonce*}[2][]{\par\medskip\noindent\textbf{#2.}\itshape}{\par\medskip}
\newenvironment{proofc}[1][\proofname]{\begin{proof}[#1]}{\end{proof}}
\newcommand{\Mc}{\mathcal{M}}
\newcommand{\pointir}{.\ ---\ }
\newcommand{\equalenv}[2]{\expandafter\let\csname#1\expandafter\endcsname\csname#2\endcsname\expandafter\let\csname end#1\expandafter\endcsname\csname end#2\endcsname}
\providecommand{\diagbox}[3][]{\begin{tikzpicture}[x=1em,y=1em,baseline=(current bounding box.center)]\useasboundingbox(0,0) rectangle(3,1.4);\draw(0,1.4)--(3,0);\node at(.5,.3){#2};\node at(2.5,1.1){#3};\end{tikzpicture}}
\providecommand{\eg}{e.g.\ }
\providecommand{\ie}{i.e.\ }
\providecommand{\cf}{cf.\ }
\providecommand{\longthanks}[1]{\par\smallskip\noindent #1\par\smallskip}
\providecommand{\qbinom}[2]{\genfrac{[}{]}{0pt}{}{#1}{#2}}

\providecommand{\mathof}[1]{\mathrm{#1}}

\numberwithin{equation}{section}


\def\endappendix{}
\let\fg\undefined

\usepackage{amsmath,amssymb,amsthm}
\usepackage{bm}
\usepackage[all]{xy}
\mathtoolsset{showonlyrefs=true}

\DeclarePairedDelimiter\abs{\lvert}{\rvert}
\newcommand{\Z}{\mathbb{Z}}
\newcommand{\N}{\mathbb{N}}
\newcommand{\Q}{\mathbb{Q}}
\newcommand{\R}{\mathbb{R}}
\newcommand{\C}{\mathbb{C}}
\newcommand{\F}{\mathbb{F}}

\newcommand{\bA}{\mathbb{A}}
\newcommand{\bB}{\mathbb{B}}
\newcommand{\bC}{\mathbb{C}}
\newcommand{\bD}{\mathbb{D}}
\newcommand{\bE}{\mathbb{E}}
\newcommand{\bF}{\mathbb{F}}
\newcommand{\bG}{\mathbb{G}}
\newcommand{\bH}{\mathbb{H}}
\newcommand{\bI}{\mathbb{I}}
\newcommand{\bJ}{\mathbb{J}}
\newcommand{\bK}{\mathbb{K}}
\newcommand{\bL}{\mathbb{L}}
\newcommand{\bM}{\mathbb{M}}
\newcommand{\bN}{\mathbb{N}}
\newcommand{\bO}{\mathbb{O}}
\newcommand{\bP}{\mathbb{P}}
\newcommand{\bQ}{\mathbb{Q}}
\newcommand{\bR}{\mathbb{R}}
\newcommand{\bS}{\mathbb{S}}
\newcommand{\bT}{\mathbb{T}}
\newcommand{\bU}{\mathbb{U}}
\newcommand{\bV}{\mathbb{V}}
\newcommand{\bW}{\mathbb{W}}
\newcommand{\bX}{\mathbb{X}}
\newcommand{\bY}{\mathbb{Y}}
\newcommand{\bZ}{\mathbb{Z}}

\newcommand{\fA}{\mathfrak{A}}
\newcommand{\fB}{\mathfrak{B}}
\newcommand{\fC}{\mathfrak{C}}
\newcommand{\fD}{\mathfrak{D}}
\newcommand{\fE}{\mathfrak{E}}
\newcommand{\fF}{\mathfrak{F}}
\newcommand{\fG}{\mathfrak{G}}
\newcommand{\fH}{\mathfrak{H}}
\newcommand{\fI}{\mathfrak{I}}
\newcommand{\fJ}{\mathfrak{J}}
\newcommand{\fK}{\mathfrak{K}}
\newcommand{\fL}{\mathfrak{L}}
\newcommand{\fM}{\mathfrak{M}}
\newcommand{\fN}{\mathfrak{N}}
\newcommand{\fO}{\mathfrak{O}}
\newcommand{\fP}{\mathfrak{P}}
\newcommand{\fQ}{\mathfrak{Q}}
\newcommand{\fR}{\mathfrak{R}}
\newcommand{\fS}{\mathfrak{S}}
\newcommand{\fT}{\mathfrak{T}}
\newcommand{\fU}{\mathfrak{U}}
\newcommand{\fV}{\mathfrak{V}}
\newcommand{\fW}{\mathfrak{W}}
\newcommand{\fX}{\mathfrak{X}}
\newcommand{\fY}{\mathfrak{Y}}
\newcommand{\fZ}{\mathfrak{Z}}

\newcommand{\fa}{\mathfrak{a}}
\newcommand{\fb}{\mathfrak{b}}
\newcommand{\fc}{\mathfrak{c}}
\newcommand{\fd}{\mathfrak{d}}
\newcommand{\fe}{\mathfrak{e}}
\newcommand{\ff}{\mathfrak{f}}
\newcommand{\fg}{\mathfrak{g}}
\newcommand{\fh}{\mathfrak{h}}
\newcommand{\fj}{\mathfrak{j}}
\newcommand{\fk}{\mathfrak{k}}
\newcommand{\fl}{\mathfrak{l}}
\newcommand{\fm}{\mathfrak{m}}
\newcommand{\fn}{\mathfrak{n}}
\newcommand{\fo}{\mathfrak{o}}
\newcommand{\fp}{\mathfrak{p}}
\newcommand{\fq}{\mathfrak{q}}
\newcommand{\fr}{\mathfrak{r}}
\newcommand{\fs}{\mathfrak{s}}
\newcommand{\ft}{\mathfrak{t}}
\newcommand{\fu}{\mathfrak{u}}
\newcommand{\fv}{\mathfrak{v}}
\newcommand{\fw}{\mathfrak{w}}
\newcommand{\fx}{\mathfrak{x}}
\newcommand{\fy}{\mathfrak{y}}
\newcommand{\fz}{\mathfrak{z}}

\newcommand{\cA}{\mathcal{A}}
\newcommand{\cB}{\mathcal{B}}
\newcommand{\cC}{\mathcal{C}}
\newcommand{\cD}{\mathcal{D}}
\newcommand{\cE}{\mathcal{E}}
\newcommand{\cF}{\mathcal{F}}
\newcommand{\cG}{\mathcal{G}}
\newcommand{\cH}{\mathcal{H}}
\newcommand{\cI}{\mathcal{I}}
\newcommand{\cJ}{\mathcal{J}}
\newcommand{\cK}{\mathcal{K}}
\newcommand{\cL}{\mathcal{L}}
\newcommand{\cM}{\mathcal{M}}
\newcommand{\cN}{\mathcal{N}}
\newcommand{\cO}{\mathcal{O}}
\newcommand{\cP}{\mathcal{P}}
\newcommand{\cQ}{\mathcal{Q}}
\newcommand{\cR}{\mathcal{R}}
\newcommand{\cS}{\mathcal{S}}
\newcommand{\cT}{\mathcal{T}}
\newcommand{\cU}{\mathcal{U}}
\newcommand{\cV}{\mathcal{V}}
\newcommand{\cW}{\mathcal{W}}
\newcommand{\cX}{\mathcal{X}}
\newcommand{\cY}{\mathcal{Y}}
\newcommand{\cZ}{\mathcal{Z}}


\newcommand{\wtil}[1]{\widetilde{#1}}
\newcommand{\ol}[1]{\overline{#1}}
\newcommand{\ul}[1]{\underline{#1}}
\newcommand{\parenth}[1]{\left( #1 \right)}
\newcommand{\lrang}[1]{\langle #1 \rangle}
\newcommand{\MEMO}[1]{\textcolor{red}{#1}}

\newcommand{\hZ}{\hat{\Z}}
\newcommand{\GamLambda}{{}_{\Gamma}\Lambda}

\newcommand{\Gal}{{\rm Gal}}
\newcommand{\Coker}{{\rm Cok}}
\newcommand{\Cok}{{\rm Cok}}
\newcommand{\Ker}{{\rm Ker}}
\newcommand{\Imag}{{\rm Im}}
\newcommand{\Div}{{\rm Div}}
\newcommand{\Pic}{{\rm Pic}}
\newcommand{\Jac}{{\rm Jac}}
\newcommand{\aug}{{\rm aug}}
\newcommand{\Frac}{{\rm Frac}}
\newcommand{\depth}{{\rm depth}}
\newcommand{\Image}{{\rm Im}}
\newcommand{\Tor}{{\rm Tor}}

\newcommand{\Fitt}{{\rm Fitt}}
\newcommand{\Ann}{{\rm Ann}}
\newcommand{\pd}{{\rm pd}}

\newcommand{\bad}{{\rm bad}}
\newcommand{\id}{{\rm id}}
\newcommand{\Cl}{{\rm Cl}}
\newcommand{\ord}{{\rm ord}}
\newcommand{\Hom}{{\rm Hom}}
\newcommand{\Ext}{{\rm Ext}}
\newcommand{\End}{{\rm End}}
\newcommand{\Ind}{{\rm Ind}}

\newcommand{\rank}{{\rm rank}}
\newcommand{\cha}{{\rm char}}













\begin{document}
\section*{\raggedright 842: Explicit norm-quotient Fitting ideals of Jacobian groups for finite connected abelian Galois graph covers}
\noindent\textbf{Author:} Kataoka, Takenori\par\noindent\textbf{Source:} Fitting ideals of Jacobian groups of graphs. Algebraic Combinatorics 7 (2024), publisher version, DOI 10.5802/alco.355. \url{https://alco.centre-mersenne.org/articles/10.5802/alco.355/}
\par\noindent\textbf{License:} CC BY 4.0, \url{https://creativecommons.org/licenses/by/4.0/}. The original publisher PDF cover has the explicit license notice. See the saved source/license records.
\par\noindent\textbf{Editorial material note (not author text):} The selected problem is Theorem1.1 for exact finite connected abelian Galois graph coverings and the authored norm quotient, cyclic decomposition and generator formulas. Full Sections1--4 retain the complete human norm/Jacobian exact sequences, graph/arithmetic reduction, augmentation-ideal tensor presentation, tree-supported maximal-minor classification, all prescribed monomial realization and terminal generator calculation. The complete original AppendixA supplies Fitting/once-shift definitions and named prior shift-theory arguments. Standard Kirchhoff and cited Fitting/shift results remain references; the independent profinite-covering, Iwasawa and self-duality programs are unselected and keep exact original printed section pointers. All diagrams are native editable xymatrix. No theorem is broadened to the whole Jacobian or nonabelian covers, and no proof is shortened or generated. Unused Sha/Cyrillic symbol declarations are omitted from the rendering preamble only; no retained body uses them. All proof text and formulas below are editable author TeX. Mathematical wording, language and proof order are retained. Only the article class, title matter and bibliography presentation are adapted for independent compilation; no proof has been generated or repaired.
\par\noindent\textbf{Editorial difficulty:} H3: The proof reduces the graph Jacobian norm quotient through an exact sequence and shifted Fitting ideals, then computes an augmentation-ideal presentation by tensor resolutions. Its crucial maximal-minor calculation classifies surviving minors by trees and derives all monomial generators, retaining the finite abelian covering and norm-quotient restrictions.
\medskip\hrule\medskip
\noindent\textbf{Author statement, complete proof and local supporting text follow in source order.}
\newtheorem{claim}[theo]{Claim}
\newtheorem{ass}[theo]{Assumption}
\section{Introduction}\label{sec:intro}

For a (finite) graph $X$, let $\Jac(X)$ denote its Jacobian group, which is also known as the sandpile group.
It is known that $\Jac(X)$ is a finite abelian group.
See \S\ref{sec:defn_J_graph} for basic notions about graphs and their Jacobian groups.

Let $Y/X$ be an abelian covering of connected graphs and $\Gamma$ its Galois group.
Then the Jacobian group $\Jac(Y)$ is naturally equipped with a $\Z[\Gamma]$-module structure, which is the main theme of the present paper.
More concretely, we focus on its Fitting ideal (see \S\ref{ss:Fitt_defn} for the definition).
The main result of this paper gives a complete description of the Fitting ideal
\[
\Fitt_{\Z[\Gamma]/(N_{\Gamma})}(\Jac(Y)/N_{\Gamma}\Jac(Y)).
\]
Here, $N_{\Gamma}$ denotes the norm element of $\Gamma$.
See \S\ref{ss:main_result} for the concrete statement.
This result can be regarded as an analogue of a result of Atsuta and the author~\cite{AK21} in number theory, as explained in Remark ~\ref{rem:AK} below.

We will also study infinite abelian coverings in a sense in \S6.
In this case, we obtain a complete description of the Fitting ideal of the associated Jacobian group itself, without dividing by the Jacobian group of the base graph.

The study of infinite coverings of graphs is morally an analogue of Iwasawa theory in number theory.
Recently such an analogue of Iwasawa theory for graphs is developed (e.g., Gonet~\cite{Gon22} and a series of work of Valli\`{e}res including~\cite{MV21b}).
The present work is directly inspired by this situation surrounding graph theory and Iwasawa theory.
Moreover, in \S8, as a digression, we give concise proofs of the analogues of the Iwasawa class number formula and of Kida's formula.
Those results are already proved by others, but the author hopes that this makes Iwasawa theory for graphs more accessible for those who are familiar with techniques in Iwasawa theory.

\subsection{Main result for finite coverings}\label{ss:main_result}

Now we state the main result.
Let $Y/X$ be an abelian covering of connected graphs and $\Gamma$ its Galois group.
In \S\ref{sec:defn_J_graph}, we will introduce an element
$
Z_{Y/X} \in \Z[\Gamma],
$
which is related to the (equivariant) Ihara zeta function via the three-term determinant formula (see \S5.3).

Let us take a decomposition
\[
\Gamma = \Delta_1 \times \cdots \times \Delta_s,
\]
where $\Delta_l$ is a cyclic group of order $n_l \geq 1$ for each $1 \leq l \leq s$.
For each $l$, we fix a generator $\sigma_l$ of $\Delta_l$ and define $\tau_l$, $\nu_l \in \Z[\Delta_l] \subset \Z[\Gamma]$ by
\[
\tau_l = \sigma_l - 1,
\quad
\nu_l = 1 + \sigma_l + \sigma_l^2 + \cdots + \sigma_l^{n_l-1}.
\]
Then $\nu_l$ is the norm element of $\Delta_l$.
The norm element of $\Gamma$, defined by $N_{\Gamma} = \sum_{\gamma \in \Gamma} \gamma$, satisfies $N_{\Gamma} = \nu_1 \cdots \nu_s$.
We also define an element $D_l \in \Z[\Delta_l]$ by
\[
D_l = \sigma_l + 2 \sigma_l^2 + \cdots + (n_l-1) \sigma_l^{n_l-1}.
\]

The main result is the following.
The proof will be given in \S\S\ref{sec:Fin_arith}--\ref{sec:fin_alg}.

\begin{thm}\label{thm:main_fin}
We have
\begin{align}
& \Fitt_{\Z[\Gamma]/(N_{\Gamma})}(\Jac(Y)/N_{\Gamma}\Jac(Y))\\
& \quad = Z_{Y/X} \sum_{l=1}^s \left( \nu_1 \cdots \nu_{l-1} \cdot \frac{D_l}{n_l} \cdot \nu_{l+1} \cdots \nu_s \right) \\
& \quad \quad + Z_{Y/X} \left(\nu_1^{e_1} \cdots \nu_s^{e_s} \tau_1^{f_1} \cdots \tau_s^{f_s} \, \middle| \, \begin{matrix} 0 \leq e_l \leq 1, f_l \geq 0, \\ e_1 + \cdots + e_s + f_1 + \cdots + f_s = s- 2 \end{matrix} \right).
\end{align}
Here the right hand side, which is literally an ideal of $\Z[\Gamma]$, is regarded as an ideal of $\Z[\Gamma]/(N_{\Gamma})$ via the natural projection.
\end{thm}

\begin{rema}\label{rem:Kol}
In the theory of Euler systems in number theory, the element $D_l$ is often called the Kolyvagin derivative operator and plays an important role.
It seems to be interesting that this element is useful in computation that is not directly related to the theory of Euler systems.
\end{rema}

In Proposition ~\ref{prop:NormY}, we will show an isomorphism
$
\Jac(X) \simeq N_{\Gamma} \Jac(Y).
$
Therefore, Theorem ~\ref{thm:main_fin} can be viewed as a description of the Fitting ideal of $\Jac(Y)/\Jac(X)$.
Unfortunately, determining $\Fitt_{\Z[\Gamma]}(\Jac(Y))$ itself seems to be out of our reach.

\begin{rema}\label{rem:AK}
Theorem ~\ref{thm:main_fin} can be regarded as an analogue of a result of Atsuta and the author~\cite{AK21} on the Fitting ideal of class groups in number theory.
Let us roughly review the result.
Let $L/K$ be a finite abelian extension of number fields such that $K$ is totally real and $L$ is a CM-field.
We consider the $T$-ray class group $\Cl_L^T$ of $L$, where $T$ is an auxiliary finite set of finite primes of $K$.
Then in~\cite{AK21}, (assuming the relevant equivariant Tamagawa number conjecture), we gave a complete description of the Fitting ideal
\[
\Fitt_{\Z[\Gal(L/K)]^-}(\Cl_L^{T, -}),
\]
where the superscript $(-)^-$ denotes the minus component with respect to the complex conjugation.
The plus component seems to be out of our reach.
This is analogous to the situation in Theorem ~\ref{thm:main_fin}.
\end{rema}

\subsection{Idea of proof}\label{ss:comments}

In number theory, there are already a number of studies on the Fitting ideals of arithmetic modules.
In particular, in the field of Iwasawa theory, the Fitting ideals of Iwasawa modules have been studied by Greither--Kurihara (e.g.,~\cite{GK15},~\cite{GKT20} with Tokio, etc.).

In~\cite{Kata_05}, the author introduced the ``shift theory'' of Fitting ideals, which we will review in \S\ref{ss:shift_defn}.
It is a useful technique to compute the Fitting ideals of arithmetic modules.
Roughly speaking, the technique has two steps: 
\begin{itemize}
\item[(1)]
we study the arithmetic module and obtain a description of the Fitting ideal using an explicit algebraic factor, and 
\item[(2)]
we compute the algebraic factor in a purely algebraic way.
\end{itemize}
The first applications of this theory in~\cite{Kata_05} were to Iwasawa theory, i.e., infinite extensions of number fields.
Then Atsuta and the author~\cite{AK21} made an application to class groups of number fields as explained in Remark ~\ref{rem:AK}.

In this paper, we again employ the shift theory as the basic technique to prove both Theorem ~\ref{thm:main_fin} and the result for infinite coverings.
Therefore, we follow the two steps explained above.
See Theorems ~\ref{thm:Fitt_fin1} and ~\ref{thm:Fitt_fin} for the first and second steps, respectively.

The algebraic factor to be computed in Theorem ~\ref{thm:Fitt_fin} seems to be a new one as discussed in Remark ~\ref{rem:Kol}.
Moreover, the proof of the formula is not easy, and it is the most technical part in this paper.
We employ a combinatoric method (using graph theory), though the problem is purely algebraic.
Note that Greither--Kurihara--Tokio~\cite{GKT20} also applied graph theory to compute Fitting ideals.

On the other hand, the algebraic factor to be computed for the infinite covering case does not seem to be very new.
Indeed, it can be computed by using a formula that has been obtained in~\cite{AK21}.
However, a relatively minor issue is that in this paper we deal with commutative rings that are not necessarily noetherian.
Therefore, we have to generalize the shift theory, which was developed only over noetherian rings in~\cite{Kata_05}.
This is explained in \S\ref{ss:shift_defn}.

\subsection{Organization of this paper}\label{ss:org}

In \S\ref{sec:defn_J_graph}, we introduce basic notions about graphs.
Then in \S\ref{sec:Fin_arith}, we reduce the proof of Theorem ~\ref{thm:main_fin} to an algebraic problem.
The algebraic problem is solved in \S\ref{sec:fin_alg}.

After preliminaries in \S5 on voltage graphs and their derived graphs, in \S6 we state and prove the result for infinite coverings.

In \S7, we observe the self-duality of the Jacobian groups, which is a contrast to the analogue in number theory.
In \S8, we explain the module-theoretic approach to Iwasawa theory for graphs.
Both \S\S7 and ~8 can be read independently.
Finally in \S\ref{sec:Fitt_gen}, we review the definition of the Fitting ideals and the shift theory.

\section{Graphs and their Jacobian groups}\label{sec:defn_J_graph}

In this section, we introduce basic notions about graphs and their Jacobian groups.
References are Corry--Perkinson~\cite[Chapters 1 and 2]{CP18} and Baker--Norine~\cite[\S 1.3]{BN07}.
One can also consult recent papers of Valli\`eres such as Hammer--Mattman--Sands--Valli\`eres~\cite[\S 2]{HMSV24} and Ray--Valli\`eres~\cite[\S 2]{RV22}.

In this paper, graphs are always assumed to be finite.
We allow graphs to have multi-edges and loops (so one may call them multigraphs).
More precisely, we define graphs as follows, using Serre's formalism~\cite[Chapter I, \S 2.1]{Ser80}.

\begin{defi}
A graph $X$ consists of a finite set $V_X$ of vertices, a finite set $\bE_X$ of edges, an automorphism of $\bE_X$ denoted by $e \mapsto \ol{e}$, and two maps $s_X$, $t_X$ from $\bE_X$ to $V_X$ (often abbreviated to $s = s_X$ and $t = t_X$), satisfying the following:
\begin{itemize}
\item
For any $e \in \bE_X$, we have $\ol{e} \neq e$ and $\ol{\ol{e}} = e$, that is, the automorphism $e \mapsto \ol{e}$ is an involution of $\bE_X$ without fixed points.
\item
For any $e \in \bE_X$, we have $s(\ol{e}) = t(e)$ and $t(\ol{e}) = s(e)$.
\end{itemize}

Each element $e \in \bE_X$ is regarded as an edge from $s(e)$ to $t(e)$, and $\ol{e}$ is regarded as the inverse of $e$.
Let us write $E_X$ for the quotient set of $\bE_X$ obtained by identifying $e$ and $\ol{e}$.
Then we have a canonical two-to-one map from $\bE_X$ to $E_X$, so their cardinalities satisfy $\# E_X = \frac{1}{2} \cdot \# \bE_X$.

For each $v \in V_X$, we define
\[
\bE_{X, v} = \{ e \in \bE_X \mid s(e) = v\},
\]
which is the set of edges that start from $v$.
\end{defi}

This formalism involving $\bE_X$ and the involution will be useful for introducing voltage graphs in \S5.
On the other hand, we can regard $X$ as an undirected \mbox{(multi-)}graph through the quotient set $E_X$.
The notions that we will introduce in this section are essentially defined for the associated undirected graph structure.


In this paper, for simplicity, we usually deal with connected graphs.
Being connected is not an essential assumption since the general case can be easily deduced.

In the rest of this section, let $X$ be a connected graph.

\begin{defi}
We define the divisor group $\Div(X)$ as the free $\Z$-module on the set~$V_X$, namely, 
\[
\Div(X) = \bigoplus_{v \in V_X} \Z [v].
\]
Let us write
\[
\deg_X: \Div(X) \to \Z
\]
for the $\Z$-homomorphism that sends $[v]$ to $1$ for any $v \in V_X$.
We define $\Div^0(X)$ as the kernel of $\deg_X$.
\end{defi}

We obviously have an exact sequence
\begin{equation}\label{eq:D0}
0 \to \Div^0(X) \to \Div(X) \overset{\deg_X}{\to} \Z \to 0.
\end{equation}

\begin{defi}\label{defn:LAD}
We define $\Z$-homomorphisms
\[
\cL_X, A_X, D_X: \Div(X) \to \Div(X)
\]
by
\[
A_X([v]) = \sum_{e \in \bE_{X, v}} [t(e)],
\quad
D_X([v]) = (\# \bE_{X, v}) [v]
\]
for $v \in V_X$, and $\cL_X = D_X - A_X$.
The presentation matrices of $\cL_X$, $A_X$, and $D_X$ (with respect to the basis $\{ [v] \}_{v \in V_X}$) are called the Laplacian matrix, the adjacency matrix, and the degree matrix, respectively.
\end{defi}

It is easy to see $\deg_X \circ \cL_X = 0$, that is, the image of $\cL_X$ is contained in $\Div^0(X)$.
Therefore, we can make the following definition.

\begin{defi}\label{defn:JP}
We define the Jacobian group $\Jac(X)$ and the Picard group $\Pic(X)$ as the cokernels
\[
\Jac(X) = \Coker(\cL_X: \Div(X) \to \Div^0(X))
\]
and
\[
\Pic(X) = \Coker(\cL_X: \Div(X) \to \Div(X)).
\]
\end{defi}

Therefore, by \eqref{eq:D0}, we have an exact sequence
\begin{equation}\label{eq:JP_ex}
0 \to \Jac(X) \to \Pic(X) \overset{\deg_X}{\to} \Z\to 0.
\end{equation}

A fundamental property of the Jacobian group is the following.

\begin{thm}[Kirchhoff's matrix tree theorem]\label{thm:Kirch}
The Jacobian group $\Jac(X)$ is a finite $\Z$-module.
In fact, the cardinality of $\Jac(X)$ is equal to the number of spanning trees of $X$.
\end{thm}

Let us observe the following homological description of the Jacobian group.

\begin{lemm}\label{lem:cpx_J}
Let $\iota_X: \Z \to \Div(X)$ be the $\Z$-homomorphism that sends $1$ to $\sum_{v \in V_X} [v]$.
Then the sequence
\begin{equation}\label{eq:cpx_J}
0 \to \Z \overset{\iota_X}{\to} \Div(X) \overset{\cL_X}{\to} \Div(X) \overset{\deg_X}{\to} \Z \to 0
\end{equation}
is a complex that is acyclic except for the right $\Div(X)$, at which the homology group is isomorphic to $\Jac(X)$.
\end{lemm}

\begin{proof}
We only have to check $\Image(\iota_X) = \Ker(\cL_X)$; the other assertions are clear.
It is easy to see that $\cL_X \circ \iota_X = 0$, that is, $\Image(\iota_X) \subset \Ker(\cL_X)$.
Thanks to Theorem ~\ref{thm:Kirch}, the $\Z$-ranks of $\Image(\iota_X)$ and $ \Ker(\cL_X)$ are the same.
Moreover, the definition of $\iota_X$ implies that $\Image(\iota_X)$ is a saturated submodule of $\Div(X)$ (i.e., the cokernel of $\iota_X$ is $\Z$-torsion-free).
These observations imply $\Image(\iota_X) = \Ker(\cL_X)$.
\end{proof}


Now we quickly recall the notion of Galois coverings of connected graphs (see Terras~\cite[Chapter 13]{Ter11} or~\cite[\S 2.1]{RV22} for the details).
A covering $Y/X$ of connected graphs means that we are given a covering map $\pi: Y \to X$, which is surjective and locally isomorphic.
The degree of $Y/X$ is defined as $\# (\pi^{-1}(v))$ for any choice of $v \in V_X$.
The Galois group $\Gamma$ of $Y/X$ is defined as the automorphism group of $Y$ that respects the covering map.
We say that a covering $Y/X$ is Galois (or normal) if the order $\# \Gamma$ of $\Gamma$ is equal to the degree of $Y/X$.
We say $Y/X$ is abelian if $\Gamma$ is abelian.


\begin{defi}\label{defn:Z2}
Let $Y/X$ be a Galois covering of connected graphs with Galois group~$\Gamma$.
It is easy to see that $\Div(Y)$ is a free $\Z[\Gamma]$-module and the endomorphism $\cL_Y$ on $\Div(Y)$ is a $\Z[\Gamma]$-homomorphism.
In case $\Gamma$ is abelian, we define
\begin{equation}\label{eq:Z_YX}
Z_{Y/X} = \det_{\Z[\Gamma]}(\cL_Y \mid \Div(Y)) \in \Z[\Gamma].
\end{equation}
\end{defi}

By the definitions of $\Pic(Y)$ and of Fitting ideals, we have
\begin{equation}\label{eq:Z_YX2}
\Fitt_{\Z[\Gamma]}(\Pic(Y)) = (Z_{Y/X})
\end{equation}
when $\Gamma$ is abelian.
In \S5.3, we will see a relation between $Z_{Y/X}$ and the Ihara zeta function.

\section{Reduction to an algebraic problem}\label{sec:Fin_arith}

In this section, we describe the Fitting ideal that is concerned in Theorem ~\ref{thm:main_fin} by using a shifted Fitting ideal of an explicit module.

We begin with an elementary lemma.
Let $\Gamma$ be a finite group and write
\[
R = \Z[\Gamma],
\qquad
\ol{R} = R/(N_{\Gamma}),
\]
where we set $N_{\Gamma} = \sum_{\gamma \in \Gamma} \gamma \in \Z[\Gamma]$.

\begin{lemm}\label{lem:Tor1}
We have
\[
\Tor_1^R(\ol{R}, \Z) = 0,
\quad
\ol{R} \otimes_R \Z \simeq \Z/(\# \Gamma)\Z.
\]
\end{lemm}

\begin{proof}
We have an exact sequence of $R$-modules
\[
0 \to \Z \overset{\iota}{\to} R \to \ol{R} \to 0,
\]
where $\iota$ sends $1$ to $N_{\Gamma}$.
Then the associated long exact sequence of $\Tor_{*}^R(-, \Z)$ reads
\[
0 \to \Tor_1^R(\ol{R}, \Z) \to \Z \otimes_R \Z \overset{\iota_*}{\to} R \otimes_R \Z \to \ol{R} \otimes_R \Z \to 0.
\]
We have natural isomorphisms
\[
\Z \otimes_R \Z \simeq \Z,
\quad
R \otimes_R \Z \simeq \Z,
\]
through which $\iota_*$ is identified with the multiplication by $\# \Gamma$ on $\Z$.
Therefore, we obtain the lemma.
\end{proof}

In the rest of this section, let $Y/X$ be a Galois covering of connected graphs.
We write $\Gamma$ for its Galois group and use the same notation $R$ and $\ol{R}$ as above.

The following is the key ingredient to prove Theorem ~\ref{thm:main_fin}.

\begin{prop}\label{prop:Jac_ex1}
We have an exact sequence of finite $\ol{R}$-modules
\[
0 \to \ol{R} \otimes_R \Jac(Y) \to \ol{R} \otimes_R \Pic(Y) \to \Z/(\# \Gamma)\Z \to 0.
\]
\end{prop}

\begin{proof}
Let us consider the exact sequence \eqref{eq:JP_ex} for $Y$ instead of $X$.
By taking $\Tor_*^R(\ol{R}, -)$, we obtain an exact sequence
\[
\Tor_1^R(\ol{R}, \Z) \to \ol{R} \otimes_R \Jac(Y) \to \ol{R} \otimes_R \Pic(Y) \to \ol{R} \otimes_R \Z \to 0.
\]
By applying Lemma ~\ref{lem:Tor1}, we obtain the proposition.
\end{proof}

Let us study $\ol{R} \otimes_R \Pic(Y)$.

\begin{lemm}\label{lem:olPic}
We have a short exact sequence
\[
0 \to\ol{R} \otimes_R \Div(Y) \overset{\ol{\cL_{Y}}}{\to} \ol{R} \otimes_R \Div(Y) \to \ol{R} \otimes_R \Pic(Y) \to 0,
\]
where $\ol{\cL_Y}$ denotes the homomorphism induced by $\cL_Y$.
\end{lemm}

\begin{proof}
By the definition of $\Pic(Y)$, the cokernel of $\ol{\cL_Y}$ is isomorphic to $\ol{R} \otimes_R \Pic(Y)$, which is finite by Theorem ~\ref{thm:Kirch} and Proposition ~\ref{prop:Jac_ex1}.
Then the injectivity of $\ol{\cL_Y}$ follows since it is an endomorphism of a free $\Z$-module of finite rank.
\end{proof}

We write $\ol{Z_{Y/X}} \in \ol{R}$ for the image of $Z_{Y/X} \in R$.
By applying the shift theory (see~\S\ref{ss:shift_defn}), we obtain the following.

\begin{thm}\label{thm:Fitt_fin1}
Suppose that $Y/X$ is an abelian covering.
Then we have
\[
\Fitt_{\ol{R}}(\ol{R} \otimes_R \Jac(Y))
= \ol{Z_{Y/X}} \Fitt^{[1]}_{\ol{R}}(\Z/ (\#\Gamma) \Z)
\]
as ideals of $\ol{R}$.
\end{thm}

\begin{proof}
As in \S\ref{ss:shift_defn}, let $\cP_{\ol{R}}$ be the category of finite $\ol{R}$-modules whose projective dimensions are $\leq 1$.
By Lemma ~\ref{lem:olPic}, we see that $\ol{R} \otimes_R \Pic(Y)$ is in $\cP_{\ol{R}}$.
Moreover, \eqref{eq:Z_YX2} implies
\[
\Fitt_{\ol{R}}(\ol{R} \otimes_R \Pic(Y)) 
= (\ol{Z_{Y/X}}).
\]
Therefore, the theorem follows from the definition of $\Fitt^{[1]}_{\ol{R}}(-)$ and Proposition ~\ref{prop:Jac_ex1}.
\end{proof}

We will determine $\Fitt^{[1]}_{\ol{R}}(\Z/ (\#\Gamma) \Z)$ in \S\ref{sec:fin_alg}, which would complete the proof of Theorem ~\ref{thm:main_fin}.

Before closing this section, we show a proposition that provides an interpretation of $\ol{R} \otimes_R \Jac(Y)$.
In the course of the proof, we reproduce the exact sequence in Proposition ~\ref{prop:Jac_ex1}.

\begin{prop}\label{prop:NormY}
We have a natural isomorphism
\[
\Jac(X) \simeq N_{\Gamma} \Jac(Y),
\]
so
\[
\ol{R} \otimes_R \Jac(Y) 
\simeq \Jac(Y)/\Jac(X).
\]
\end{prop}

\begin{proof}
We make use of Lemma ~\ref{lem:cpx_J} for both $X$ and $Y$.
Let us consider the following diagram involving those complexes.
\[
\xymatrix{
	0 \ar[r]
	& \Z \ar[r]^-{\iota_X} \ar@{=}[d]
	& \Div(X) \ar[r]^{\cL_X} \ar@{^(->}[d]
	& \Div(X) \ar[r]^-{\deg_X} \ar@{^(->}[d]
	& \Z \ar[r] \ar@{^(->}[d]^{\times (\# \Gamma)}
	& 0\\
	0 \ar[r]
	& \Z \ar[r]_-{\iota_Y} \ar[d]
	& \Div(Y) \ar[r]_{\cL_{Y}} \ar@{->>}[d]
	& \Div(Y) \ar[r]_-{\deg_{Y}} \ar@{->>}[d]
	& \Z \ar[r] \ar@{->>}[d]
	& 0\\
	& 0 \ar[r] 
	& \ol{R} \otimes_{R} \Div(Y) \ar[r]_{\ol{\cL_{Y}}}
	& \ol{R} \otimes_{R} \Div(Y) \ar[r]
	& \Z/(\# \Gamma) \Z \ar[r] 
	& 0.
}
\]
Here, the injective homomorphism $\Div(X) \hookrightarrow \Div(Y)$ is defined by sending $[v]$ to $\sum_{w \in \pi^{-1}(v)} [w]$, where $\pi: Y \to X$ is the covering map.
Then the image of this homomorphism is $N_{\Gamma} \Div(Y)$.
The vertical sequences are all exact and the lower horizontal sequence is the induced complex.

We regard this large diagram as a short exact sequence of complexes.
Recall that the kernel and the cokernel of $\ol{\cL_Y}$ are determined in Lemma ~\ref{lem:olPic}.
Then, taking the homology groups, we obtain an exact sequence
\begin{equation}\label{eq:31-1}
0 \to \Jac(X) \to \Jac(Y) \to \ol{R} \otimes_R \Pic(Y)  \to \Z/(\# \Gamma) \Z \to 0.
\end{equation}
The construction of the embedding of $\Jac(X)$ into $\Jac(Y)$ shows that the image coincides with $N_{\Gamma} \Jac(Y)$, as claimed.
Moreover, this reproduces the exact sequence in Proposition ~\ref{prop:Jac_ex1}.
\end{proof}

\section{Algebraic result}\label{sec:fin_alg}

Let $\Gamma$ be any finite abelian group.
Set $R = \Z[\Gamma]$ and $\ol{R} = \Z[\Gamma]/(N_{\Gamma})$.
The goal of this section is to determine the factor $\Fitt^{[1]}_{\ol{R}}(\Z/ (\#\Gamma) \Z)$ that appeared in Theorem ~\ref{thm:Fitt_fin1}.

\subsection{Statement}\label{ss:fin_alg}

First we state the result.
As in \S\ref{ss:main_result}, we fix a decomposition $\Gamma = \Delta_1 \times \cdots \times \Delta_s$ into cyclic groups.
For each $1 \leq l \leq s$, we choose a generator $\sigma_l$ of $\Delta_l$ and define $n_l = \# \Delta_l$ and elements $\tau_l, \nu_l, D_l \in \Z[\Delta_l]$ by the same formulas.
Clearly we have $\tau_l \nu_l = 0$.
A key property of $D_l$ is 
\begin{equation}\label{eq:D_key}
\tau_l D_l = n_l - \nu_l,
\end{equation}
which can be proved by a direct computation.
This property is also important in the theory of Euler systems.

The result is the following.

\begin{thm}\label{thm:Fitt_fin}
We have
\begin{align}
\Fitt^{[1]}_{\ol{R}}(\Z/ (\#\Gamma) \Z)
= & \sum_{l=1}^s \left( \nu_1 \cdots \nu_{l-1} \cdot \frac{D_l}{n_l} \cdot \nu_{l+1} \cdots \nu_s \right) \\
& + \left(\nu_1^{e_1} \cdots \nu_s^{e_s} \tau_1^{f_1} \cdots \tau_s^{f_s} \, \middle| \, \begin{matrix} 0 \leq e_l \leq 1, f_l \geq 0, \\ e_1 + \cdots + e_s + f_1 + \cdots + f_s = s- 2 \end{matrix} \right).
\end{align}
\end{thm}

Then clearly Theorem ~\ref{thm:main_fin} follows from Theorems ~\ref{thm:Fitt_fin1} and ~\ref{thm:Fitt_fin}.
The rest of this section is devoted to the proof of Theorem ~\ref{thm:Fitt_fin}.

First we describe $\Fitt^{[1]}_{\ol{R}}(\Z/ (\#\Gamma) \Z)$ by using an (unshifted) Fitting ideal.
Let 
\[
I = \Ker(\Z[\Gamma] \to \Z) \subset R
\]
 be the augmentation ideal, which we can regard as an $\ol{R}$-module.

\begin{lemm}\label{lem:Fitt1}
We have
\[
\Fitt^{[1]}_{\ol{R}}(\Z/ (\#\Gamma) \Z) 
 = (\# \Gamma)^{-1} \Fitt_{\ol{R}}(I/R(\# \Gamma - N_{\Gamma})).
\]
\end{lemm}

\begin{proof}
By the tautological exact sequence $0 \to I \to R \to \Z \to 0$, we have an exact sequence $0 \to I \to \ol{R} \to \Z/(\# \Gamma)\Z \to 0$.
Since $\# \Gamma - N_{\Gamma}$ is in $I$ and its image to $\ol{R}$ is $\# \Gamma$, we then obtain an exact sequence
\[
0 \to I/R (\# \Gamma - N_{\Gamma}) \to \ol{R}/ (\# \Gamma) \ol{R} \to \Z/(\# \Gamma)\Z \to 0.
\]
By the definition of $\Fitt_{\ol{R}}^{[1]}(-)$, this implies the lemma.
\end{proof}

By Lemma ~\ref{lem:Fitt1}, the proof of Theorem ~\ref{thm:Fitt_fin} is reduced to the following.

\begin{prop}\label{prop:Fitt_fin2}
We have
\begin{align}
& \Fitt_{R}(I/R(\# \Gamma - N_{\Gamma}))\\
& \qquad = 
	(N_{\Gamma})
	+ (\# \Gamma) \cdot \sum_{l=1}^s \left( \nu_1 \cdots \nu_{l-1} \cdot \frac{D_l}{n_l} \cdot \nu_{l+1} \cdots \nu_s \right) \\
& \qquad \qquad + (\# \Gamma) \cdot \left( \nu_1^{e_1} \cdots \nu_s^{e_s} \tau_1^{f_1} \cdots \tau_s^{f_s} \, \middle| \, \begin{matrix} 0 \leq e_l \leq 1, f_l \geq 0, \\ e_1 + \cdots + e_s + f_1 + \cdots + f_s = s- 2 \end{matrix} \right).
\end{align}
\end{prop}

The proof of this proposition will be given in \S\ref{ss:thm2}.
Before that, we establish a key technical proposition in \S\ref{ss:resol}.

\subsection{Key proposition}\label{ss:resol}

Let $s \geq 0$ be an integer and let $R_s = \Z[T_1, \dots, T_s]$ be the polynomial ring in $s$ indeterminates.
Let us construct a free resolution of $\Z$ over $R_s$ by using tensor products of complexes.
This, or rather the more complicated construction in \S\ref{ss:thm2} below, is inspired by~\cite{GK15} and is also used in~\cite{Kata_05},~\cite{AK21}, etc.

For each $1 \leq l \leq s$, we have an exact sequence
\[
0 \to \Z[T_l] x_l \overset{T_l}{\to} \Z[T_l] \to \Z \to 0.
\]
Here, $x_l$ denotes an indeterminate, so $\Z[T_l] x_l$ is a free $\Z[T_l]$-module of rank one, and the map labeled $T_l$ sends $x_l$ to $T_l$.
The map $\Z[T_l] \to \Z$ sends $T_l$ to $0$.

Observe that $R_s$ is the tensor product of $\Z[T_l]$ over $\Z$ for $1 \leq l \leq s$.
Then by taking the tensor product over $\Z$ of the complexes $[\Z[T_l] x_l \overset{T_l}{\to} \Z[T_l]]$ for $1 \leq l \leq s$, we obtain an exact sequence
\[
\bigoplus_{1 \leq l < l' \leq s} R_s x_l x_{l'} \overset{d_2}{\to} \bigoplus_{1 \leq l \leq s} R_s x_l \overset{d_1}{\to} R_s \overset{d_0}{\to} \Z \to 0,
\]
where $d_0$ sends $T_l$ to $0$ for any $l$, $d_1$ is determined by 
\[
d_1(x_l) = T_l
\]
for $1 \leq l \leq s$, and $d_2$ is determined by
\[
d_2(x_l x_{l'}) = - T_{l'} x_l + T_l x_{l'}
\]
for $1 \leq l < l' \leq s$ (there are other choices of signs, but that does not matter).

We write $N_s(T_1, \dots, T_s)$ for the presentation matrix of $d_2$ with respect to the basis $\{x_l x_{l'} \mid 1 \leq l < l' \leq s \}$ and $x_1, \dots, x_s$.
Note that we do not determine the order of the rows because it does not matter at all.
Indeed, in the following we will sometimes choose various orders of bases that are suitable for computation.

\begin{exam}
When $s = 0$ (resp.~$s = 1$), by definition the source of $d_2$ is the zero module, so $d_2 = 0$ and $N_0()$ (resp.~$N_1(T_1)$) is the empty matrix.
When $s = 2$, we have 
\[
N_2(T_1, T_2) = \begin{pmatrix} -T_2 & T_1 \end{pmatrix}.
\]
When $s = 3$, we have
\[
N_3(T_1, T_2, T_3) = 
\begin{pmatrix}
& - T_3 & T_2\\
- T_3 & & T_1\\
- T_2 & T_1 &
\end{pmatrix}.
\]
Here, we use the order $x_2x_3, x_1x_3, x_1x_2$ for the basis.
When $s = 4$, we have
\[
N_4(T_1, T_2, T_3, T_4) = 
\begin{pmatrix}
-T_2 & T_1 &  & \\
- T_3 & & T_1 & \\
- T_4 & & & T_1\\
& - T_3 & T_2 & \\
& -T_4 & & T_2\\
& & -T_4 & T_3
\end{pmatrix}.
\]
Here, we use the order $x_1 x_2, x_1x_3, x_1x_4, x_2x_3, x_2x_4, x_3x_4$ for the basis.
\end{exam}

For a matrix $H$, which is identified with a homomorphism between free modules, we have the associated ideal $\Fitt(H)$ defined as in Definition ~\ref{defn:Fitt1}.

The following formula plays a key role in the proof of Proposition ~\ref{prop:Fitt_fin2}.
It is the most technical result in this paper.

\begin{prop}\label{prop:Fitt_p2}
We have
\[
\Fitt(\begin{pmatrix} N_s(T_1, \dots, T_s) \\ \begin{matrix} B_1 & \cdots & B_s \end{matrix} \end{pmatrix})
= \begin{cases}
(1) & \text{if $s = 0$,}\\
(B_1) & \text{if $s = 1$,}\\
\left( \sum_{l=1}^s B_l T_l \right)(T_1, \dots, T_s)^{s - 2} & \text{otherwise.}
\end{cases}
\]
Here, $B_1, \dots, B_s$ are indeterminates that have no relation with $T_1, \dots, T_s$, i.e., we work over the polynomial ring $\Z[T_1, \dots, T_s, B_1, \dots, B_s]$.
\end{prop}

The rest of this subsection is devoted to the proof of this proposition.
The cases $s = 0$, $s = 1$ are clear.
Let us suppose $s \geq 2$.

Recall that the rows of $N_s(T_1, \dots, T_s)$ are labeled $x_l x_{l'}$ ($1 \leq l < l' \leq s$).
We also attach a label $y$ to the last row $\begin{pmatrix} B_1 & \cdots & B_s \end{pmatrix}$.
For any subset 
\[
\cA \subset \{x_l x_{l'} \mid 1 \leq l < l' \leq s\} \amalg \{y\}
\]
with $\# \cA = s$, let $N_{\cA}$ be the $s \times s$ submatrix that is constructed by picking up the rows whose labels are in $\cA$.
Then the left hand side in Proposition ~\ref{prop:Fitt_p2} is generated by $\det(N_{\cA})$ for all such $\cA$.

For such an $\cA$, let us construct an undirected simple graph $G_{\cA}$ that has $s$ vertices $x_1, \dots, x_s$ so that $x_l x_{l'} \in \cA$ if and only if $x_l$ and $x_{l'}$ are adjacent.
Then the number of the edges of $G_{\cA}$ is either $s-1$ or $s$; indeed, it is $s-1$ if and only if $y \in \cA$.
This construction gives a one-to-one correspondence between the set
\[
\{ \cA \subset \{x_l x_{l'} \mid 1 \leq l < l' \leq s\} \amalg \{y\} \mid \# \cA = s\}
\]
and the set
\[
\{ \text{ simple graph structures on the set of vertices $\{x_1, \dots, x_s\}$ with $s-1$ or $s$ edges }\}.
\]

We shall describe $\det(N_{\cA})$ by using information about the associated graph $G_{\cA}$.

\begin{claim}\label{claim2}
We have $\det(N_{\cA}) = 0$ unless $G_{\cA}$ is a tree.
\end{claim}

\begin{proof}
Suppose that $G_{\cA}$ is not a tree.
Since the number of the edges of $G_{\cA}$ is at least $s-1$, this assumption is equivalent to that $G_{\cA}$ has a cycle.
In other words, by permutation of the indices, we may assume that $\{x_1x_2, x_2 x_3, \dots, x_{r-1} x_r, x_r x_1\} \subset \cA$ for some $3 \leq r \leq s$.
In this case, the matrix $N_{\cA}$ is of the form
\[
\begin{pmatrix}
	\begin{matrix}
		-T_2 & T_1 & & & \\
		& -T_3 & T_2 & & \\
		& & \ddots & \ddots & \\
		& & & - T_r & T_{r-1}\\
		-T_r & & & & T_1
	\end{matrix}
 & 0\\
* & * \\
\end{pmatrix}
\]
if we use the order
\[
x_1, x_2, \dots, x_r, *, \dots, *
\]
of the vertices and the order 
\[
x_1x_2, x_2 x_3, \dots, x_{r-1} x_r, x_r x_1, *, \dots, *
\]
 of the rows in $\cA$ ($*$ denotes unspecified things).
It is easy to see that the $r \times r$ matrix in the upper left has determinant $0$.
Therefore, the claim follows.
\end{proof}

By Claim ~\ref{claim2}, we only have to deal with the case where $G_{\cA}$ is a tree.
Note that then the number of the edges of $G_{\cA}$ is $s-1$, i.e., $y \in \cA$.
For each $1 \leq l \leq s$, let $\deg_{\cA}(x_l) \geq 1$ be the degree of $x_l$ in the graph $G_{\cA}$.
By definition, $\deg_{\cA}(x_l)$ is the number of vertices that are adjacent to $x_l$ in $G_{\cA}$.

\begin{claim}\label{claim3}
If $G_{\cA}$ is a tree, then we have
\[
\det(N_{\cA}) = \pm \left(\sum_{l = 1}^s B_l T_l \right) \cdot T_1^{\deg_{\cA}(x_1)-1} \cdots T_s^{\deg_{\cA}(x_s) - 1}.
\]
\end{claim}

\begin{proof}
Let us write simply $\deg(-)$ instead of $\deg_{\cA}(-)$.
First we show that the claim follows from the following weaker claim: there exist signs $\epsilon_1, \dots, \epsilon_s \in \{\pm 1\}$ such that we have
\[
\det(N_{\cA}) = \left(\sum_{l = 1}^s \epsilon_l B_l T_l \right) \cdot T_1^{\deg(x_1)-1} \cdots T_s^{\deg(x_s)-1}.
\]
Suppose that such a family $\{\epsilon_l\}_l$ exists.
By Claim ~\ref{claim2}, for any $l \neq l'$, we know that $\det(N_{\cA})$ vanishes if we set $B_l = T_{l'}$, $B_{l'} = -T_l$, and the other $B$'s to be zero.
Therefore, $\epsilon_l T_{l'} T_l + \epsilon_{l'} (-T_l) T_{l'} = 0$, which shows $\epsilon_l = \epsilon_{l'}$.
Thus we obtain $\epsilon_1 = \cdots = \epsilon_s$, so the full claim follows.

Let us show the weaker claim.
By the cofactor expansion of $\det(N_{\cA})$ with respect to the final row (labeled $y$), we have
\[
\det(N_{\cA})
= \sum_{l = 1}^s B_l \delta_l,
\]
where $\delta_l$ denotes the cofactor of $B_l$ (i.e., the determinant of the submatrix obtained by eliminating the last row $y$ and the $l$-th column).
Here and henceforth, we ignore the sign of $\delta_l$, which does not matter for the weaker claim.
Then it is enough to show that, for any fixed $l$, the cofactor $\delta_l$ coincides with $T_l \cdot T_1^{\deg(x_1)-1} \cdots T_s^{\deg(x_s)-1}$ up to sign.

Fix $l$.
Let us reorder the vertices $x_1, \dots, x_s$ as follows.
We regard $x_l$ as the root of the tree $G_{\cA}$.
Recall that then the depth of each vertex $x_k$, denoted by $\depth(x_k)$, is defined as the length of the unique path from the root $x_l$ to $x_k$.
(We set $\depth(x_l) = 0$.)
Now we reorder the vertices of the rooted graph $G_{\cA}$ as
\[
x_{\sigma(1)}, x_{\sigma(2)}, \dots, x_{\sigma(s)},
\]
where $\sigma$ is a permutation of the set $\{1, 2, \dots, s\}$ such that 
\[
\depth(x_{\sigma(k)}) \leq  \depth(x_{\sigma(k+1)})
\]
for every $1 \leq k \leq s - 1$.
We necessarily have $\sigma(1) = l$, but this $\sigma$ is not unique in general.

For each $2 \leq k \leq s$, let $1 \leq \ul{k} \leq s$ be the index such that $x_{\sigma(\ul{k})}$ is the parent of $x_{\sigma(k)}$.
It is the unique vertex that is adjacent to $x_{\sigma(k)}$ and whose depth is less than that of $x_{\sigma(k)}$.
Note that we have $1 \leq \ul{k} < k$.

Now, to compute the cofactor $\delta_l$, we use the order
\[
x_{\sigma(2)}, x_{\sigma(3)}, \dots, x_{\sigma(s)}
\]
of the vertices and the order
\[
x_{\sigma(\ul{2})} x_{\sigma(2)}, x_{\sigma(\ul{3})} x_{\sigma(3)}, \dots, x_{\sigma(\ul{s})} x_{\sigma(s)}
\]
of the edges.
Then, thanks to $\ul{k} < k$, the matrix whose determinant is $\delta_l$ is lower triangular and we obtain
\[
\delta_l = \pm T_{\sigma(\ul{2})} T_{\sigma(\ul{3})} \cdots T_{\sigma(\ul{s})}.
\]
Since $x_{\sigma(\ul{k})}$ is the parent of $x_{\sigma(k)}$, for each $1 \leq l' \leq s$, the exponent of the indeterminate $T_{l'}$ in this product is equal to the number of the children of the vertex $x_{l'}$.
If the vertex $x_{l'}$ is the root, i.e., if $l' = l$, then the number of children is equal to $\deg(x_l)$.
Otherwise, i.e., if $l' \neq l$, the number of children is equal to $\deg(x_{l'}) - 1$.
This completes the proof of the claim.
\end{proof}

\begin{exam}
We illustrate the above proof by an example.
Let $s = 5$ and consider $\cA = \{x_1 x_2, x_1 x_3, x_1 x_4, x_2 x_5, y\}$, so
\[
N_{\cA}
= \begin{pmatrix}
	-T_2 & T_1 & & &\\
	-T_3 & & T_1 & &\\
	-T_4 & & & T_1 &\\
	& -T_5 & & & T_2\\
	B_1 & B_2 & B_3 & B_4 & B_5	
\end{pmatrix}.
\]
In this case we have $\deg_{\cA}(x_1) = 3$, $\deg_{\cA}(x_2) = 2$, and $\deg_{\cA}(x_3) = \deg_{\cA}(x_4) = \deg_{\cA}(x_5) = 1$.

To consider the cofactor $\delta_1$ of $B_1$, we regard $x_1$ as the root.
Then $\depth(x_2) = \depth(x_3) = \depth(x_4) = 1$ and $\depth(x_5) = 2$, so we use the order $x_2$, $x_3$, $x_4$, $x_5$ for the vertices.
We have $\ul{2} = 1$, $\ul{3} = 1$, $\ul{4} = 1$, $\ul{5} = 2$, so we use the order $x_1 x_2$, $x_1 x_3$, $x_1 x_4$, $x_2 x_5$ for the edges.
Then we obtain
\[
\delta_1 = \pm \det 
\begin{pmatrix}
	T_1 & & &\\
	& T_1 & &\\
	& & T_1 &\\
	T_5 & & & T_2\\
\end{pmatrix}
= \pm T_1^3 T_2.
\]
Here and in the following examples, we omit writing $\pm$ before the indeterminates since we may ignore the signs.

Similarly, to consider $\delta_2$, we use $x_1$, $x_5$, $x_3$, $x_4$ for the vertices and $x_1 x_2$, $x_2 x_5$, $x_1 x_3$, $x_1 x_4$ for the edges, and obtain
\[
\delta_2 = \pm \det 
\begin{pmatrix}
	T_2 & & &\\
	& T_2 & &\\
	T_3 & & T_1 &\\
	T_4 & & & T_1\\
\end{pmatrix}
= \pm T_1^2 T_2^2.
\]

To consider $\delta_3$, we use the orders $x_1$, $x_2$, $x_4$, $x_5$ and $x_1 x_3$, $x_1 x_2$, $x_1 x_4$, $x_2 x_5$ and obtain
\[
\delta_3 = \pm \det 
\begin{pmatrix}
	T_3 & & &\\
	T_2 & T_1 & &\\
	T_4 & & T_1 &\\
	& T_5 & & T_2\\
\end{pmatrix}
= \pm T_1^2 T_2 T_3.
\]
\end{exam}

Now we return to the general case.
By Claims ~\ref{claim2} and ~\ref{claim3}, the Fitting ideal to be computed in Proposition ~\ref{prop:Fitt_p2} is generated by
\[
\left(\sum_{l = 1}^s B_l T_l \right) \cdot T_1^{\deg_{\cA}(x_1)-1} \cdots T_s^{\deg_{\cA}(x_s) - 1},
\]
where $G_{\cA}$ varies all tree structures on the set of vertices $\{x_1, \dots, x_s \}$.
Note that we have
\[
\sum_{l=1}^s (\deg_{\cA}(x_l) - 1) 
= \sum_{l=1}^s \deg_{\cA}(x_l) - s
= 2(s-1) - s = s-2.
\]
It remains only to show that, conversely, the tuple $(\deg_{\cA}(x_1) - 1, \dots, \deg_{\cA}(x_s) - 1)$ can be any tuple whose sum is $s - 2$.
This assertion follows from the following.

\begin{claim}
Let $s \geq 2$ and let $f_1, \dots, f_s$ be integers such that $f_l \geq 0$ and $f_1 + \dots + f_s = s - 2$.
Then there exists a tree structure on the set of vertices $x_1, \dots, x_s$ such that the degree of $x_l$ equals $f_l + 1$ for all $1 \leq l \leq s$.
\end{claim}

\begin{proof}
We argue by the induction on $s$.
When $s = 2$, we must have $f_1 = f_2 = 0$, and the unique tree structure satisfies the property.
Suppose $s \geq 3$.
By $f_1 + \dots + f_s = s - 2$, we have $f_l = 0$ and $f_{l'} \geq 1$ for some $l$ and $l'$, so we may assume that $f_s = 0$ and $f_1 \geq 1$.
By the induction hypothesis, there exists a tree structure on $x_1, \dots, x_{s-1}$ such that the degree of $x_1$ is $f_1$ and the degree of $x_l$ is $f_l+1$ for $2 \leq l \leq s-1$.
Then we obtain the desired graph by simply connecting $x_1$ and $x_s$.
\end{proof}

This completes the proof of Proposition ~\ref{prop:Fitt_p2}.

\subsection{Proof of Proposition ~\ref{prop:Fitt_fin2}}\label{ss:thm2}

In this subsection, we prove Proposition ~\ref{prop:Fitt_fin2}.

We first construct a free resolution of the augmentation ideal $I \subset R = \Z[\Gamma]$ over $R$ in a similar way as in \S\ref{ss:resol}.
Note that the idea is already used in previous work such as~\cite{AK21}.

For each $1 \leq l \leq s$, we have an exact sequence
\[
\Z[\Delta_l] x_l^2 \overset{\nu_l}{\to} \Z[\Delta_l] x_l \overset{\tau_l}{\to} \Z[\Delta_l] \to \Z \to 0,
\]
where $x_l$ denotes an indeterminate, the map $\nu_l$ sends $x_l^2$ to $\nu_l x_l$, $\tau_l$ sends $x_l$ to $\tau_l$, and the map $\Z[\Delta_l] \to \Z$ is the augmentation map.

By taking the tensor product of the complexes $[\Z[\Delta_l] x_l^2 \overset{\nu_l}{\to} \Z[\Delta_l] x_l \overset{\tau_l}{\to} \Z[\Delta_l]]$ over $\Z$, we obtain an exact sequence
\[
\bigoplus_{l=1}^s R x_l^2 \oplus \bigoplus_{1 \leq l < l' \leq s} R x_l x_{l'} 
\overset{d_2}{\to} \bigoplus_{l=1}^s R x_l 
\overset{d_1}{\to} R
\overset{d_0}{\to} \Z \to 0,
\]
where $d_0$ is the augmentation map, $d_1$ is determined by
\[
d_1(x_l) = \tau_l  \quad (1 \leq l \leq s),
\]
and $d_2$ is determined by
\[
\begin{cases}
d_2(x_l^2) = \nu_l x_l & (1 \leq l \leq s),\\
d_2(x_l x_{l'}) = -\tau_{l'} x_l + \tau_l x_{l'} &(1 \leq l < l' \leq s).
\end{cases}
\]

Therefore, the presentation matrix of $d_2$ is of the form
\[
M_s(\nu_1,\dots, \nu_s, \tau_1, \dots, \tau_s) = 
\begin{pmatrix}
\begin{matrix}
\nu_1 & & \\
& \ddots & \\
& & \nu_s \\
\end{matrix}\\
 N_s(\tau_1, \dots, \tau_s) 
\end{pmatrix}.
\]
Here, $N_s(\tau_1, \dots, \tau_s)$ denotes the matrix obtained by setting $T_l = \tau_l$ in the matrix $N_s(T_1, \dots, T_s)$ constructed in \S\ref{ss:resol}.

We have constructed a presentation $M_s(\nu_1,\dots, \nu_s, \tau_1, \dots, \tau_s)$ of the $\Z[\Gamma]$-module $I = \Ker(d_0: R \to \Z)$.
Our next task is to construct a presentation of $I/R(\# \Gamma - N_{\Gamma})$.
For that purpose, we define elements $b_1, \cdots, b_s \in R$ by
\[
b_l = \nu_1 \cdots \nu_{l-1} \cdot D_l \cdot n_{l+1} \cdots n_s
\]
for $1 \leq l \leq s$.
This is where the Kolyvagin derivative operators come into play.

\begin{lemm}\label{lem:sum_bt}
We have
\[
\sum_{l=1}^s b_l \tau_l = \# \Gamma - N_{\Gamma}.
\]
\end{lemm}

\begin{proof}
By using the identity \eqref{eq:D_key}, we can compute
\begin{align}
\sum_{l=1}^s b_l \tau_l 
& = \sum_{l=1}^s \nu_1 \cdots \nu_{l-1} \cdot (n_l - \nu_l) \cdot n_{l+1} \cdots n_s\\
& = n_1 \cdots n_s - \nu_1 \cdots \nu_s\\
& = \# \Gamma - N_{\Gamma}.
\end{align}
Thus we obtain the lemma.
\end{proof}

Then Lemma ~\ref{lem:sum_bt} implies that
\[
d_1 \left(\sum_{l=1}^s b_l x_l \right) = \# \Gamma - N_{\Gamma},
\]
so $I/R(\# \Gamma - N_{\Gamma})$ has a presentation matrix
\[
\begin{pmatrix} M_s(\nu_1, \dots, \nu_s, \tau_1, \dots, \tau_s) \\ \begin{matrix} b_1 & \cdots & b_s \end{matrix} \end{pmatrix}
 = \begin{pmatrix} \begin{matrix} \nu_1 & & \\
& \ddots & \\
& & \nu_s \\
\end{matrix}\\
 N_s(\tau_1, \dots, \tau_s) \\
 \begin{matrix} b_1 & \cdots & b_s \end{matrix}
 \end{pmatrix}
\]
over $R$.

To  compute the Fitting ideal of this matrix,
we first observe
\begin{align}\label{eq:tM0}
& \Fitt(
	\begin{pmatrix} 
		M_s(\nu_1, \dots, \nu_s, \tau_1, \dots, \tau_s) \\ \begin{matrix} b_1 & \cdots & b_s \end{matrix}
	\end{pmatrix}
) \\
& \quad = \sum_{j = 0}^s \sum_{\substack{{\bm a} \subset \{1, 2, \dots, s\} \\ \# {\bm a} = j}} \nu_{a_1} \cdots \nu_{a_j} 
\Fitt(
	\begin{pmatrix} N_{s - j}(\tau_{a_{j + 1}}, \dots, \tau_{a_s}) \\ \begin{matrix} b_{a_{j+1}} & \cdots & b_{a_s} \end{matrix}
	\end{pmatrix}
).
\end{align}
Here, $\bm{a}$ runs over the subsets of $\{1, 2, \dots, s\}$ with $\# \bm{a} = j$ and define $a_1, \dots, a_s$ by requiring 
\[
\bm{a} = \{a_1, \dots, a_j\},
\quad \{a_1, \dots, a_s\} = \{1, 2, \dots, s\}, 
\quad a_1 < \dots < a_j, 
\quad a_{j + 1} < \dots < a_s.
\]
We can show \eqref{eq:tM0} directly by using the identity $\nu_l \tau_l = 0$ (see~\cite[Proposition 4.7]{AK21} for a similar reasoning).

Let us compute the right hand side of \eqref{eq:tM0} by using Proposition ~\ref{prop:Fitt_p2}.
When $j = s$, we have $\bm{a} = \{1, 2, \dots, s\}$, so the term is 
\begin{equation}\label{eq:tM1}
\nu_1 \cdots \nu_s (1) = (N_{\Gamma}).
\end{equation}
When $j = s - 1$, we have $\{1, 2, \dots, s\} \setminus \bm{a} = \{ l \}$ for some $1 \leq l \leq s$, and then the term is
\begin{equation}\label{eq:tM2}
\nu_1 \cdots \nu_{l-1} \cdot \nu_{l+1} \cdots \nu_s \cdot (b_l)
= (\# \Gamma) \cdot \left( \nu_1 \cdots \nu_{l-1} \cdot \frac{D_l}{n_l} \cdot \nu_{l+1} \cdots \nu_s \right).
\end{equation}
Finally we consider $0 \leq j \leq s - 2$.
For each $\bm{a}$ with $\# \bm{a} = j$, the term can be computed as
\begin{align}\label{eq:tM3}
& \nu_{a_1} \cdots \nu_{a_j} \left( b_{a_{j+1}} \tau_{a_{j+1}} + \cdots + b_{a_s} \tau_{a_s} \right)(\tau_{a_{j+1}}, \dots, \tau_{a_s})^{s - j - 2}\\
& \quad = \nu_{a_1} \cdots \nu_{a_j} \left( b_{a_1} \tau_{a_1} + \cdots + b_{a_s} \tau_{a_s} \right)(\tau_{a_{1}}, \dots, \tau_{a_s})^{s - j - 2}\\
& \quad = \nu_{a_1} \cdots \nu_{a_j} ( \# \Gamma - N_{\Gamma}) (\tau_1, \dots, \tau_s)^{s - j - 2},
\end{align}
where the first equality follows from the identity $\tau_l \nu_l = 0$; the second from Lemma ~\ref{lem:sum_bt}.

It is easy to see that the ideal generated by \eqref{eq:tM1}, \eqref{eq:tM2}, and \eqref{eq:tM3} coincides with the right hand side of the proposition (observe that, thanks to \eqref{eq:tM1}, the $N_{\Gamma}$ in \eqref{eq:tM3} can be ignored).
This completes the proof of Proposition ~\ref{prop:Fitt_fin2}.

This also completes the proof of Theorem ~\ref{thm:main_fin}.
\appendix

\section{Notes on Fitting ideals}\label{sec:Fitt_gen}

In this section, we briefly review the definition of Fitting ideals and the shift theory introduced in~\cite{Kata_05}.

Let $R$ be a commutative ring.
Note that we do not assume $R$ is noetherian.
This is because the coefficient rings arising from the arithmetic in this paper are not noetherian in general (e.g., $\hZ[[\hZ]]$).
In~\cite{Kata_05}, the author developed the shift theory only over noetherian rings.
In this section, we observe that the argument can be generalized to non-noetherian rings by imposing appropriate finiteness properties on modules.

\subsection{Fitting ideals}\label{ss:Fitt_defn}

We recall the definition of Fitting ideals.
See, e.g., Northcott~\cite[\S 3.1]{Nor76}.

\begin{defi}\label{defn:Fitt1}
Let $I$ be a finite set and $J$ a (not necessarily finite) set.
Let
\[
h: R^{\oplus J} \to R^{\oplus I}
\]
be an $R$-homomorphism, where $R^{\oplus I}$, $R^{\oplus J}$ denote the free modules on the set $I$, $J$ respectively.

We define the (initial) Fitting ideal $\Fitt(h) = \Fitt_R(h)$ as follows.
For each subset $J' \subset J$ with $\# J' = \# I$, we write 
\[
h_{J'}: R^{\oplus J'} \to R^{\oplus I}
\]
for the restriction of $h$.
Let $\det(h_{J'})$ be the determinant of $h_{J'}$ with respect to any choice of bases; inevitably this determinant has ambiguity up to $R^{\times}$, but this does not matter.
Then $\Fitt_R(h)$ is defined as the ideal of $R$ generated by $\det(h_{J'})$ for all $J' \subset J$ with $\# J' = \# I$.
Note that if $\# J < \# I$, we have $\Fitt_R(h) = 0$ as there is not such a~$J'$.

More generally, for an integer $i \geq 0$, the $i$-th Fitting ideal $\Fitt_i(h) = \Fitt_{i, R}(h)$ is defined as follows.
For subsets $I' \subset I$ and $J' \subset J$ with $\# I' = \# J'$, we write 
\[
h_{J', I'}: R^{\oplus J'} \to R^{\oplus I'}
\]
for the map induced by $h$.
If $i \leq \# I$, we define $\Fitt_{i, R}(h)$ as the ideal of $R$ generated by $\det(h_{J', I'})$ for all $J' \subset J$, $I' \subset I$ with $\# J' = \# I' = \# I - i$.
If $i > \# I$, we set $\Fitt_{i, R}(h) = R$.
\end{defi}

\begin{defi}\label{defn:Fitt2}
Let $M$ be a finitely generated $R$-module.
We define $\Fitt_R(M)$ and $\Fitt_{i, R}(M)$ ($i \geq 0$) respectively as $\Fitt_R(h)$ and $\Fitt_{i, R}(h)$, where $h$ is a homomorphism as in Definition ~\ref{defn:Fitt1} such that the cokernel of $h$ is isomorphic to $M$.
It is known that this definition is independent from the choice of $h$.
\end{defi}

\subsection{Shifts of Fitting ideals}\label{ss:shift_defn}

We write $\Frac(R)$ for the total ring of fractions of~$R$.
An $R$-module $M$ is said to be torsion if any element is annihilated by a non-zero-divisor of $R$; equivalently if $\Frac(R) \otimes_R M = 0$.

\begin{lemm}\label{lem:tors_gen}
Let $F$ be a free $R$-module of finite rank and $h: F \to F$ be an endomorphism.
Then the following are equivalent.
\begin{itemize}
\item[(i)]
The map $h$ is injective.
\item[(ii)]
The determinant $\det(h) \in R$ is a non-zero-divisor.
\item[(iii)]
The cokernel $\Coker(h)$ is a torsion $R$-module.
\end{itemize}
\end{lemm}

\begin{proof}
The fact (i) $\Leftrightarrow$ (ii) is shown in~\cite[Chapter III, \S 8, Proposition 3, p.~524]{BouAlg74}.
Since $\Cok(h)$ is annihilated by $\det(h)$, we have (ii) $\Rightarrow$ (iii).
If (iii) holds, the base change of $h$ from $R$ to $\Frac(R)$ is surjective, so it is isomorphic, which implies (i).
\end{proof}

Let us define $\cP_R$ as the category of finitely presented torsion $R$-module $P$ such that $\pd_R(P) \leq 1$, where $\pd_R(-)$ denotes the projective dimension.
By Lemma ~\ref{lem:tors_gen}, if a module $P$ satisfies an exact sequence of the form
\begin{equation}\label{eq:pres}
0 \to F \to F \to P \to 0
\end{equation}
with $F$ a free module of finite rank, then $P$ is in $\cP_R$.
Conversely, if $R$ is local, then every $P$ in $\cP_R$ satisfies an exact sequence of the form \eqref{eq:pres}, since in that case any projective module is necessarily free.

A fractional ideal $I$ of $R$ is defined as an $R$-submodule of $\Frac(R)$ such that $u I \subset R$ for some non-zero-divisor $u \in R$.

\begin{lemm}\label{lem:FittP}
The following are true.
\begin{itemize}
\item[(1)]
For each $P \in \cP_R$, the Fitting ideal $\Fitt_R(P)$ is invertible as a fractional ideal of $R$.
\item[(2)]
Let $0 \to M' \to M \to P \to 0$ be an exact sequence of finitely generated torsion $R$-modules such that $P \in \cP_R$.
Then we have
\[
\Fitt_R(M) = \Fitt_R(P) \Fitt_R(M').
\]
\end{itemize}
\end{lemm}

\begin{proof}
We sketch the proof (see~\cite[Proposition 2.7]{Kata_05} for the details).
For both claims (1) and (2), it is enough to show them after localization at all prime ideals of $R$, so we may assume $R$ is local.
Then $P$ satisfies an exact sequence of the form \eqref{eq:pres}.
Now claim (1) follows from Lemma ~\ref{lem:tors_gen} and claim (2) follows by using the horseshoe lemma.
\end{proof}

Now we introduce the ``$n$-th shift'' $\Fitt^{[n]}_R(-)$ as in~\cite[Theorem 2.6]{Kata_05}.
First we consider the ``once-shift'' $\Fitt^{[1]}_R(-)$, which actually suffices for the purpose of this paper.
For completeness, in Definition ~\ref{defn:shift_n}, we will also introduce the general $n$-th shift.

\begin{defi}
For a finitely presented torsion $R$-module $M$, we define a fractional ideal $\Fitt^{[1]}_R(M)$ as follows.
Let us take an exact sequence of $R$-modules
\begin{equation}\label{eq:resol1}
0 \to N \to P \to M \to 0
\end{equation}
with $P \in \cP_R$.
For instance, if $M$ is generated by $n$ elements and annihilated by a non-zero-divisor $f \in R$, we may take $P = (R/fR)^n$.
Since $M$ is finitely presented, $N$ is finitely generated, so its Fitting ideal is defined.
Then we define (using Lemma ~\ref{lem:FittP}(1))
\[
\Fitt^{[1]}_R(M) = \Fitt_R(P)^{-1} \Fitt_R(N).
\]
This is well-defined, that is, independent from the choice of \eqref{eq:resol1}.
\end{defi}

Let us sketch the proof of the independency from \eqref{eq:resol1}.
Let $0 \to N' \to P' \to M \to 0$ be another exact sequence with $P' \in \cP_R$.
Then, by using the pull-back $L$ of the maps $P \to M$ and $P' \to M$, we obtain a commutative diagram with exact rows and columns
\[
\xymatrix{
	& & N' \ar@{=}[r] \ar@{^(->}[d]
	& N' \ar@{^(->}[d]
	& \\
	0 \ar[r]
	& N \ar[r] \ar@{=}[d]
	& L \ar[r] \ar@{->>}[d]
	& P' \ar[r] \ar@{->>}[d]
	& 0\\
	0 \ar[r]
	& N \ar[r]
	& P \ar[r]
	& M \ar[r]
	& 0.
}
\]
Since $P, P' \in \cP_R$, by Lemma ~\ref{lem:FittP}(2), we obtain
\[
\Fitt_R(L) = \Fitt_R(P') \Fitt_R(N),
\quad
\Fitt_R(L) = \Fitt_R(P) \Fitt_R(N').
\]
These formulas imply 
\[
\Fitt_R(P)^{-1} \Fitt_R(N)
= \Fitt_R(P')^{-1} \Fitt_R(N'),
\]
as desired.

Finally let us introduce the general $n$-th shift.
See~\cite[Theorem 2.6]{Kata_05} for the proof; to remove the noetherian hypothesis, we only have to suitably deal with the finiteness conditions.

\begin{defi}\label{defn:shift_n}
Let $n \geq 0$ be an integer.
Let $M$ be a torsion $R$-module such that there exists an exact sequence
\[
R^{a_n} \to R^{a_{n-1}} \to \dots \to R^{a_0} \to M \to 0
\]
for some integers $a_0, \dots, a_n \geq 0$.
For instance, when $n = 0$ (resp.~$n = 1$), this condition means that $M$ is finitely generated (resp.~finitely presented).
For such an $M$, we define a fractional ideal $\Fitt^{[n]}_R(M)$ as follows.
Let us take an exact sequence of $R$-modules
\begin{equation}\label{eq:resol2}
0 \to N \to P_1 \to \cdots \to P_n \to M \to 0
\end{equation}
with $P_1, \dots, P_n \in \cP_R$.
The existence of such a sequence follows from the assumption on $M$, and moreover then $N$ is finitely generated.
Then we define (using Lemma ~\ref{lem:FittP}(1))
\[
\Fitt^{[n]}_R(M) = \left( \prod_{i=1}^n \Fitt_R(P_i)^{(-1)^i} \right) \Fitt_R(N).
\]
This is well-defined, that is, independent from the choice of \eqref{eq:resol2}.
\end{defi}
 
\begin{thebibliography}{MMMMMM}
\bibitem[1]{AK21} Atsuta, Mahiro; Kataoka, Takenori Fitting ideals of class groups for CM abelian extensions, Algebra Number Theory, Volume 17 (2023) no. 11, pp. 1901-1924 \href{https://alco.centre-mersenne.org/articles/10.5802/alco.355/#AK21}{Publisher reference}.
\bibitem[2]{BN07} Baker, Matthew; Norine, Serguei Riemann-Roch and Abel-Jacobi theory on a finite graph, Adv. Math., Volume 215 (2007) no. 2, pp. 766-788 \href{https://alco.centre-mersenne.org/articles/10.5802/alco.355/#BN07}{Publisher reference}.
\bibitem[3]{BouAlg74} Bourbaki, Nicolas Elements of mathematics. Algebra, Part I: Chapters 1-3, Hermann, Paris; Addison-Wesley Publishing Co., Reading, Mass., 1974, xxiii+709 pages (translated from the French) \href{https://alco.centre-mersenne.org/articles/10.5802/alco.355/#BouAlg74}{Publisher reference}.
\bibitem[4]{CP18} Corry, Scott; Perkinson, David Divisors and sandpiles: An introduction to chip-firing, American Mathematical Society, Providence, RI, 2018, xiv+325 pages \href{https://alco.centre-mersenne.org/articles/10.5802/alco.355/#CP18}{Publisher reference}.
\bibitem[5]{Gon22} Gonet, Sophia R. Iwasawa Theory of Jacobians of Graphs, Algebr. Comb., Volume 5 (2022) no. 5, pp. 827-848 \href{https://alco.centre-mersenne.org/articles/10.5802/alco.355/#Gon22}{Publisher reference}.
\bibitem[6]{Gree99} Greenberg, Ralph Iwasawa theory for elliptic curves, Arithmetic theory of elliptic curves (Cetraro, 1997) (Lecture Notes in Math.), Volume 1716, Springer, Berlin, 1999, pp. 51-144 \href{https://alco.centre-mersenne.org/articles/10.5802/alco.355/#Gree99}{Publisher reference}.
\bibitem[7]{GK15} Greither, C.; Kurihara, M. Tate sequences and Fitting ideals of Iwasawa modules, Algebra i Analiz, Volume 27 (2015) no. 6, pp. 117-149 \href{https://alco.centre-mersenne.org/articles/10.5802/alco.355/#GK15}{Publisher reference}.
\bibitem[8]{GKT20} Greither, Cornelius; Kurihara, Masato; Tokio, Hibiki The second syzygy of the trivial G-module, and an equivariant main conjecture, Development of Iwasawa theory—the centennial of K. Iwasawa’s birth (Adv. Stud. Pure Math.), Volume 86, Math. Soc. Japan, Tokyo, 2020, pp. 317-349 \href{https://alco.centre-mersenne.org/articles/10.5802/alco.355/#GKT20}{Publisher reference}.
\bibitem[9]{HMSV24} Hammer, Kyle; Mattman, Thomas; Sands, Jonathan; Vallières, Daniel The special value u=1 of Artin-Ihara L-functions, Proc. Amer. Math. Soc., Volume 152 (2024) no. 2, pp. 501-514 \href{https://alco.centre-mersenne.org/articles/10.5802/alco.355/#HMSV24}{Publisher reference}.
\bibitem[10]{Iwa59} Iwasawa, Kenkichi On Γ-extensions of algebraic number fields, Bull. Amer. Math. Soc., Volume 65 (1959), pp. 183-226 \href{https://alco.centre-mersenne.org/articles/10.5802/alco.355/#Iwa59}{Publisher reference}.
\bibitem[11]{Kata_05} Kataoka, Takenori Fitting invariants in equivariant Iwasawa theory, Development of Iwasawa theory—the centennial of K. Iwasawa’s birth (Adv. Stud. Pure Math.), Volume 86, Math. Soc. Japan, Tokyo, 2020, pp. 413-465 \href{https://alco.centre-mersenne.org/articles/10.5802/alco.355/#Kata_05}{Publisher reference}.
\bibitem[12]{Kid80} Kida, Yûji l-extensions of CM-fields and cyclotomic invariants, J. Number Theory, Volume 12 (1980) no. 4, pp. 519-528 \href{https://alco.centre-mersenne.org/articles/10.5802/alco.355/#Kid80}{Publisher reference}.
\bibitem[13]{KM22} Kleine, Sören; Müller, Katharina On the growth of the Jacobians in ℤ p l -voltage covers of graphs, 2022 \href{https://alco.centre-mersenne.org/articles/10.5802/alco.355/#KM22}{Publisher reference}.
\bibitem[14]{MV21b} McGown, Kevin; Vallières, Daniel On abelian \(\ell\)-towers of multigraphs III, Ann. Math. Qué., Volume 48 (2024) no. 1, pp. 1-19 \href{https://alco.centre-mersenne.org/articles/10.5802/alco.355/#MV21b}{Publisher reference}.
\bibitem[15]{NSW08} Neukirch, Jürgen; Schmidt, Alexander; Wingberg, Kay Cohomology of number fields, Grundlehren der Mathematischen Wissenschaften [Fundamental Principles of Mathematical Sciences], 323, Springer-Verlag, Berlin, 2008, xvi+825 pages \href{https://alco.centre-mersenne.org/articles/10.5802/alco.355/#NSW08}{Publisher reference}.
\bibitem[16]{Nor76} Northcott, D. G. Finite free resolutions, Cambridge University Press, Cambridge-New York-Melbourne, 1976, xii+271 pages (Cambridge Tracts in Mathematics, No. 71) \href{https://alco.centre-mersenne.org/articles/10.5802/alco.355/#Nor76}{Publisher reference}.
\bibitem[17]{RV22} Ray, Anwesh; Vallières, Daniel An analogue of Kida’s formula in graph theory, 2022 \href{https://alco.centre-mersenne.org/articles/10.5802/alco.355/#RV22}{Publisher reference}.
\bibitem[18]{Ser80} Serre, Jean-Pierre Trees, Springer-Verlag, Berlin-New York, 1980, ix+142 pages (translated from the French by John Stillwell) \href{https://alco.centre-mersenne.org/articles/10.5802/alco.355/#Ser80}{Publisher reference}.
\bibitem[19]{Ter11} Terras, Audrey Zeta functions of graphs: A stroll through the garden, Cambridge Studies in Advanced Mathematics, 128, Cambridge University Press, Cambridge, 2011, xii+239 pages \href{https://alco.centre-mersenne.org/articles/10.5802/alco.355/#Ter11}{Publisher reference}.
\bibitem[20]{Was97} Washington, Lawrence C. Introduction to cyclotomic fields, Graduate Texts in Mathematics, 83, Springer-Verlag, New York, 1997, xiv+487 pages \href{https://alco.centre-mersenne.org/articles/10.5802/alco.355/#Was97}{Publisher reference}.
\end{thebibliography}
\end{document}
